\documentclass[a4paper]{article}
\usepackage[margin=1in]{geometry}
\usepackage{amsmath,amsfonts,mathtools,amsthm,amssymb}
\usepackage[l2tabu,orthodox]{nag}
\usepackage{microtype} % fixes spacing or whatever
\usepackage{enumitem}
\usepackage{tikz-cd}
\usepackage{xcolor}
\usepackage{physics}
\usepackage{mathrsfs}
\usepackage{cancel}
% \usepackage{libertine,libertinust1math}
\usepackage{eucal}
\usepackage[toc,page]{appendix}
\usepackage{pgfplots}
\pgfplotsset{compat=1.18}

\newtheorem{proposition}{Proposition}
\newtheorem{theorem}{Theorem}
\newtheorem{lemma}{Lemma}
\newtheorem{corollary}{Corollary}

\theoremstyle{definition}
\newtheorem{remark}{Remark}
\newtheorem{definition}{Definition}
\newtheorem{example}{Example}

% \usepackage[sc,osf]{mathpazo}   % With old-style figures and real smallcaps.
% \linespread{1.025}              % Palatino leads a little more leading
% % Euler for math and numbers
% \usepackage[euler-digits,small]{eulervm}
% \AtBeginDocument{\renewcommand{\hbar}{\hslash}}

\usepackage{etoolbox}
\usepackage{dynkin-diagrams}
\def\row#1/#2!{\dynkin{#1}{#2}\\}
\newcommand{\tble}[1]{
   \renewcommand*\do[1]{\row##1!}
   \[
      \begin{array}{ll}\docsvlist{#1}\end{array}
   \]
}

\allowdisplaybreaks

\renewcommand{\epsilon}{\varepsilon}
% \renewcommand{\phi}{\varphi}

\DeclareMathAlphabet{\mathpzc}{OT1}{pzc}{m}{it}

\newcommand{\goth}[1]{\mathfrak{#1}}
\newcommand{\red}[1]{#1_{\text{red}}}
\newcommand{\ad}[1]{#1_{\text{Ad}}}
\newcommand{\reg}[1]{#1_{\text{reg}}}
\newcommand{\sh}[1]{#1^{\text{sh}}}
\newcommand{\cpdv}[2]{\cfrac{\partial #1}{\partial #2}}

\newcommand{\fA}{\mathcal{A}}
\newcommand{\fB}{\mathcal{B}}
\newcommand{\fC}{\mathcal{C}}
\newcommand{\fD}{\mathcal{D}}
\newcommand{\fE}{\mathcal{E}}
\newcommand{\fF}{\mathcal{F}}
\newcommand{\fG}{\mathcal{G}}
\newcommand{\fH}{\mathcal{H}}
\newcommand{\fI}{\mathcal{I}}
\newcommand{\fJ}{\mathcal{J}}
\newcommand{\fK}{\mathcal{K}}
\newcommand{\fL}{\mathcal{L}}
\newcommand{\fM}{\mathcal{M}}
\newcommand{\fN}{\mathcal{N}}
\newcommand{\fO}{\mathcal{O}}
\newcommand{\fP}{\mathcal{P}}
\newcommand{\fQ}{\mathcal{Q}}
\newcommand{\fR}{\mathcal{R}}
\newcommand{\fS}{\mathcal{S}}
\newcommand{\fT}{\mathcal{T}}
\newcommand{\fX}{\mathcal{X}}
\newcommand{\fY}{\mathcal{Y}}
\newcommand{\fZ}{\mathcal{Z}}

\newcommand{\be}{\mathbf{e}}
\newcommand{\bx}{\mathbf{x}}
\newcommand{\bone}{\mathbf{1}}
\newcommand{\cx}{\bar{x}}
\newcommand{\cy}{\bar{y}}
\newcommand{\vx}{\vec{x}}
\newcommand{\vsigma}{\vec{\sigma}}
\newcommand{\tzeta}{\tilde{\zeta}}

\newcommand{\gm}{\goth{m}}
\newcommand{\gp}{\goth{p}}
\newcommand{\gF}{\goth{F}}
\newcommand{\gU}{\goth{U}}
\newcommand{\gV}{\goth{V}}

\newcommand{\A}{\mathbf{A}}
\newcommand{\C}{\mathbf{C}}
\newcommand{\F}{\mathbf{F}}
\newcommand{\G}{\mathbf{G}}
\newcommand{\bH}{\mathbf{H}}
\newcommand{\N}{\mathbf{N}}
\newcommand{\PP}{\mathbf{P}}
\newcommand{\Q}{\mathbf{Q}}
\newcommand{\R}{\mathbf{R}}
\newcommand{\Z}{\mathbf{Z}}

\newcommand{\dlog}{\dd\log}

\newcommand\srestr[2]{{\left.\kern-\nulldelimiterspace #1\vphantom{\small|} \right|_{#2}}}
\newcommand\restr[2]{{\left.\kern-\nulldelimiterspace #1 \vphantom{\big|} \right|_{#2}}}
\newcommand\gsec[2]{\Gamma({#1}, {#2})}
\newcommand\gsecx[1]{\Gamma(X, {#1})}

\DeclareMathOperator{\chr}{char}
\DeclareMathOperator{\rProj}{\mathbf{Proj}}
\DeclareMathOperator{\rSpec}{\mathbf{Spec}}
\DeclareMathOperator{\rHom}{\mathbf{Hom}}
\DeclareMathOperator{\rExt}{\mathbf{Ext}}
\DeclareMathOperator{\id}{id}
\DeclareMathOperator{\Frac}{Frac}
\DeclareMathOperator{\rk}{rank}
\DeclareMathOperator{\deter}{det}
\DeclareMathOperator{\pic}{Pic}
\DeclareMathOperator{\cacl}{CaCl}
\DeclareMathOperator{\trd}{tr.d.}
\DeclareMathOperator{\cl}{Cl}
\DeclareMathOperator{\depth}{depth}
\DeclareMathOperator{\codim}{codim}
\DeclareMathOperator{\Div}{Div}
\DeclareMathOperator{\coker}{coker}
\DeclareMathOperator{\len}{length}
\DeclareMathOperator{\height}{height}
\DeclareMathOperator{\supp}{Supp}
\DeclareMathOperator{\proj}{Proj}
\DeclareMathOperator{\im}{im}
\DeclareMathOperator{\Hom}{Hom}
\DeclareMathOperator{\Der}{Der}
\DeclareMathOperator{\spec}{Spec}
\DeclareMathOperator{\Aut}{Aut}
\DeclareMathOperator{\ch}{char}
\DeclareMathOperator{\tor}{Tor}
\DeclareMathOperator{\Ann}{Ann}
\DeclareMathOperator{\Syl}{Syl}
\DeclareMathOperator{\Sym}{Sym}
\DeclareMathOperator{\GL}{GL}
\DeclareMathOperator{\SL}{SL}
\DeclareMathOperator{\Stab}{Stab}
\DeclareMathOperator{\Perm}{Perm}
\DeclareMathOperator{\Orb}{Orb}
\DeclareMathOperator{\Gal}{Gal}
\DeclareMathOperator{\Supp}{Supp}
\DeclareMathOperator{\Ext}{Ext}
\DeclareMathOperator{\Tor}{Tor}
\DeclareMathOperator{\PGL}{PGL}
\DeclareMathOperator{\hd}{hd}
\DeclareMathOperator{\pd}{pd}
\DeclareMathOperator{\SO}{SO}
\DeclareMathOperator{\SU}{SU}
\DeclareMathOperator{\Sp}{Sp}
\DeclareMathOperator{\Oth}{O}
\DeclareMathOperator{\Uni}{U}
\DeclareMathOperator{\evf}{ev}
\DeclareMathOperator{\cH}{H}
\DeclareMathOperator{\dR}{R}
\DeclareMathOperator{\End}{End}
\DeclareMathOperator{\Hasse}{Hasse}
\DeclareMathOperator{\Num}{Num}

\author{James Lee}

% Solarized
% \pagecolor[RGB]{0,20,26}
% \color[RGB]{191,191,191}

% Sephia
% \pagecolor[RGB]{249,239,220}

% Grey
% \pagecolor[RGB]{200, 200, 200}

% Black
% \pagecolor[RGB]{0,0,0}
% \color[RGB]{255,255,255}

\title{Chapter 3, Section 5}

\begin{document}
\maketitle
\begin{enumerate} [label=\textbf{\arabic*.}, leftmargin=0em]

\item Let $X$ be a projective scheme over a field $k$, and let $\fF$ be a coherent sheaf on $X$.
We define the \textit{Euler characteristic} of $\fF$ by
\begin{equation*}
  \chi(\fF) = \sum (-1)^i \dim_k{\cH^i(X, \fF)}.
\end{equation*}
If
\[ \begin{tikzcd}
0 \arrow[r] & \fF' \arrow[r] & \fF \arrow[r] & \fF'' \arrow[r] & 0
\end{tikzcd} \]
is a short exact sequence of coherent sheaves on $X$, show that $\chi(\fF) = \chi(\fF') + \chi(\fF'')$.

\begin{proof}
  We get an exact sequence of finite dimensional $k$-vector spaces (5.2)
  \[ \begin{tikzcd}
    \cdots \arrow[r, "\delta_{i - 1}"] & {\cH^i(X, \fF')} \arrow[r] & {\cH^i(X, \fF)} \arrow[r] & {\cH^i(X, \fF'')} \arrow[r, "\delta_{i}"] & \cdots
  \end{tikzcd}, \]
  where the $\delta_i$'s are connecting homomorphisms.
  By exactness, we have
  \begin{align*}
    \chi(\fF) & =  \sum (-1)^i \dim_k \cH^i(X, \fF) \\
    & = \sum (-1)^i(\dim_k \coker{\delta_{i - 1}} + \dim_k \ker{\delta_i}) \\
    & = \sum (-1)^i (\dim_k \cH^i(X, \fF') - \dim_k{\im{\delta_{i -1 }}} + \dim_k \cH^i(X, \fF'') - \dim_k{\im{\delta_i}}) \\
    & = \sum (-1)^i \dim_k \cH^i(X, \fF') + \sum (-1)^i \dim_k \cH^i(X, \fF'') \\
    & = \chi(\fF') + \chi(\fF'').
  \end{align*}
\end{proof}

\item \begin{itemize}
  \item[(a)] Let $X$ be a projective scheme over a field $k$, let $\fO_X(1)$ be a very ample invertible sheaf on $X$ over $k$, and let $\fF$ be a coherent sheaf on $X$.
  Show that there is a polynomial $P(z) \in \Q[z]$, such that $\chi(\fF(n)) = P(n)$ for all $n \in \Z$. We call $P$ the \textit{Hilbert polynomial} of $\fF$ with respect to the sheaf $\fO_X(1)$.
  \item[(b)] Now let $X = \PP_k^r$, and let $M = \Gamma_*(\fF)$, considered as a graded $S = k[x_0, \dots, x_r]$-module. Use (5.2) to show that the Hilbert polynomial of $\fF$ just defined is the same as the Hilbert polynomial of $M$ defined in (I, \S 7).
\end{itemize}

\begin{proof} $ $ \vspace{0pt}
  \begin{itemize} [leftmargin=0cm]
    \item[(a)] Since $\fO_X(1)$ is a very ample sheaf on $X$ over $k$, there is a closed immersion $i : X \to \PP_k^r$ of schemes over $k$, for some $r$, such that $\fO_X(1) = i^* \fO_{\PP^r}(1)$.
    If $\fF$ is coherent on $X$, then $i_* \fF$ is coherent on $\PP_k^r$ (II, Ex. 5.5), and the cohomology is the same (2.10).
    The Euler characteristic is defined in terms of cohomology, and the projection formula (II, Ex. 5.1d) gives us
    \begin{equation*}
      i_*(\fF(n)) = i_*(\fF \otimes_{\fO_X} i^* \fO_{\PP_r}(n)) \cong i_*\fF \otimes_{\fO_{\PP^r}} \fO_{\PP^r}(n) = (i_*\fF)(n).
    \end{equation*}
    Thus, we reduce to the case $X = \PP_k^r$.

    Cover $X$ by a finite number of open affines $U_i \cong \spec{k[y_1, \dots, y_r]}$.
    Then there exist finitely generated $k[y_1, \dots, y_r]$-modules $M_i$ such that $\restr{\fF}{U_i} \cong \tilde{M}_i$.
    Since $\spec{A_i} \cap \supp{\fF}$ is a closed subset of $\spec{A_i}$; in fact, $\spec{A_i} \cap \supp{\fF} = V(\Ann{M_i})$ (AM, Ex. 3.19), $\Supp{\fF}$ is a closed subset of $X$.
    We proceed by induction on $\dim{\supp{\fF}}$.
    If $\supp{\fF}$ has dimension zero, then $\supp{\fF} \cap U_i$ is affine (II, Ex. 3.11b) equal to spectra of a Noetherian ring $A$ of dimension zero, which are precisely Artinian rings (AM, 8.5).
    So $U_i \cap \supp{\fF}$ is a finite set (AM, Ex. 8.2), and since $X$ can be covered by a finite number of $U_i$, $\Supp{\fF}$ is a finite set of closed points.
    Thus, $\fF$ is isomorphic to a direct sum of skyscraper sheaves.
    More precisely, $\fF = \bigoplus_{P \in \Supp{\fF}} \fF_P$, where $\fF_P$ is the stalk of $\fF$ at $P$, and cohomology commutes with direct sums, so by (Ex. 1), we reduce to the case when $\Supp{\fF}$ consists of a single point.
    Then, $\chi(\fF(n))$ is just the Hilbert polynomial of $\fF_P$ (I, 7.5) as a graded $A[y_1, \dots, y_r]$-module.
    
    Assume the statement to be true for all coherent sheaves with support dimension less than $n$.
    Consider the exact sequence
    \[ \begin{tikzcd}
      0 \arrow[r] & \fO_X(-1) \arrow[r, "x_r"] & \fO_X \arrow[r] & \fO_H \arrow[r] & 0,
    \end{tikzcd} \]
    where $H$ is the hyperplane $x_r = 0$.
    Taking the tensor product of this sequence with $\fF$ gives us
    \[ \begin{tikzcd}
      0 \arrow[r] & \fR \arrow[r] & \fF(-1) \arrow[r] & \fF \arrow[r] & \fF_H \arrow[r] & 0,
    \end{tikzcd} \]
    where $\fR$ is the kernel of the map $\fF(-1) \to \fF$.
    Twisting by $\fO_X(n)$ is exact, so we get
    \begin{equation*}
      \chi(\fF(n)) - \chi(\fF(n - 1)) = \chi(\fF_H(n)) - \chi(\fR(n))
    \end{equation*}
    by (III, Ex. 5.1).
    Consider the coherent $\fO_X$-module $\fR$. It vanishes under the multiplication map $x_r$, where $x_r$ locally defines one of the coordinates such that $x_r = 0$  is the defining equation for $H$, and such that there is a inclusion $\Supp{\fR} \subseteq H$, which exists by Bertini's theorem.
    Since $H \cong \PP^{r - 1}_k$, we can apply the induction hypothesis to $\fR$ and $\fF_H$, so let $P_\fR(z), P_{\fF_H}(z) \in \Q[z]$ be polynomials such that
    \begin{equation*}
      \chi(\fR(n)) = P_\fR(n), \quad \chi(\fF_H(n)) = P_{\fF_H}(n).
    \end{equation*}
    The difference function of $\chi(\fF(n))$ is equal to $P_{\fF_H}(z) - P_{\fR}(n) \in \Q[z]$, so by (\S 1, 7.3), we conclude that $\chi(\fF(n)) = P(n)$ for some $P(z) \in \Q[z]$.

    \item[(b)] By (5.2), there exists an integer $n_0$ such that for each $i > 0$ and each $n \geq n_0$, $H^i(X, \fF(n)) = 0$.
    Let $P(n)$ be the Hilbert polynomial just defined, and let $P'(n)$ be the polynomial from associated to the graded $S$-module $M$ from (I, \S 7).
    Then $P(n) = \dim_k H^0(X, \fF(n))$ for each $n \geq n_0$, but $M_n = \Gamma(X, \fF(n)) = H^0(X, \fF(n))$ and $P'(n) = \dim_k M_n$; in particular, $P(n)$ and $P'(n)$ agree at infinitely many values.
    Hence, they must be equal.
  \end{itemize}
\end{proof}

\item \textit{Arithmetic Genus.} Let $X$ be a projective scheme of dimension $r$ over a field $k$. We define the \textit{arithmetic genus $p_a$} of $X$ by
\begin{equation*}
  p_a(X) = (-1)^r (\chi(\fO_X) - 1).
\end{equation*}
Note that it depends only on $X$, not on any projective embedding.
\begin{itemize}
  \item[(a)] If $X$ is integral, and $k$ algebraically closed, show that $\cH^0(X, \fO_X) \cong k$, so that
  \begin{equation*}
    p_a(X) = \sum_{i = 0}^{r - 1} (-1)^i \dim_k{\cH^{r-i}(X, \fO_X)}.
  \end{equation*}
  In particular, if $X$ is a curve, we have
  \begin{equation*}
    p_a(X) = \dim_k{\cH^1(X, \fO_X)}.
  \end{equation*}
  \item[(b)] If $X$ is a closed subvariety of $\PP_k^r$, show that this $p_a(X)$ coincides with the one defined in (I, Ex. 7.2), which apparently depended on the projective embedding.
  \item[(c)] If $X$ is a nonsingular projective curve over an algebraically closed field $k$, show that $p_a(X)$ is in fact a \textit{birational} invariant. Conclude that a nonsingular plane curve of degree $d \geq 3$ is not rational.
\end{itemize}

\begin{proof} $ $ \vspace{0pt}
  \begin{itemize} [leftmargin=0cm]
    \item[(a)] For any embedding $j : X \to \PP_k^r$, $\cH^0(X, \fO_X) \cong k$ (I, 3.4).

    \item[(b)] Let $\fI_X$ be the ideal sheaf of $X$ on $\PP_k^r$.
    The homogenous coordinate ring $S(X)$ of $X$ for the given embedding is $k[x_0, \dots, x_r] / I$, where $I$ is the ideal $\Gamma_*(\fI_X)$ (II, Ex. 5.14).
    The result follows by (Ex. 1) and (Ex. 2b).

    \item[(c)] Two nonsingular projective curves are birational if and only if they are isomorphic (II, 6.7), so they must have the same cohomology.
    If $X$ is a nonsingular plane curve of degree $d \geq 3$, then $$p_a(X) = \frac{1}{2}(d - 1)(d - 2) > 0$$ by (I, Ex. 7.2), but $p_a(\PP^1) = 0$.
  \end{itemize}
\end{proof}

% \item[\textbf{6.}] \textit{Curves on a Nonsingular Quadric Surface.} Let $Q$ be the nonsingular quadric surface $xy = zw$ in $X = \PP_k^3$ over a field $k$. We will consider locally principal closed subschemes $Y$ of $Q$.
% These correspond to Cartier divisors on $Q$ by (II, 6.17.1).
% On the other hand, we know that $\pic{Q} \cong \Z \oplus \Z$, so we can talk about the \textit{type $(a, b)$} of $Y$ (II, 6.16) and (II, 6.6.1). Let us denote the invertible sheaf $\fL(Y)$ by $\fO_Q(a, b)$.
% Thus, for any $n \in \Z$, $\fO_Q(n) = \fO_Q(n, n)$.
% \begin{itemize}
%   \item[(a)] Use the special cases $(q, 0)$ and $(0, q)$, with $q > 0$, when $Y$ is a disjoint union of $q$ lines $\PP^1$ in $Q$, to show:
%   \begin{itemize}
%     \item[(1)] if $|a - b| \leq 1$, then $H^1(Q, \fO_Q(a, b)) = 0$;
%     \item[(2)] if $a, b < 0$, then $H^1(Q, \fO_Q(a, b)) = 0$;
%     \item[(3)] if $a \leq -2$, then $H^1(Q, \fO_Q(a, 0)) \neq 0$.
%   \end{itemize}
%   \item[(b)] Now use these results to show:
%   \begin{itemize}
%     \item[(1)] if $Y$ is a locally principal closed subscheme of type $(a, b)$ with $a, b > 0$, the $Y$ is connected;
%     \item[(2)] now assume $k$ is algebraically closed. Then for any $a, b > 0$, there exists an irreducible nonsingular curve $Y$ of type $(a, b)$. Use (II, 7.6.2) and (II, 8.18).
%     \item[(3)] an irreducible nonsingular curve $Y$ of type $(a, b)$, $a, b > 0$ on $Q$ is projectively normal (II, Ex. 5.14) if and only if $|a - b| \leq 1$. In particular, this gives lots of examples of nonsingular, but not projectively normal curves in $\PP^3$. The simplest is the one of type $(1, 3)$, which is just the rational quartic curve (I, Ex. 3.18).
%   \end{itemize}
%   \item[(c)] If $Y$ is a locally principal subscheme of type $(a, b)$ in $Q$, show that $p_a(Y) = ab - a - b + 1 = (a - 1)(b - 1)$.
% \end{itemize}

\item[\textbf{7.}]  Let $X$ (respectively, $Y$) be proper schemes over a Noetherian ring $A$.
We note by $\fL$ an invertible sheaf.
\begin{itemize}
  \item[(a)] If $\fL$ is ample on $X$, and $Y$ is any closed subscheme of $X$, then $i^* \fL$ is ample on $Y$, where $i : Y \to X$ is the inclusion.
  \item[(b)] $\fL$ is ample on $X$ if and only if $\red{\fL} = \fL \otimes \fO_{\red{X}}$ is ample on $\red{X}$.
  \item[(c)] Suppose $X$ is reduced. Then $\fL$ is ample on $X$ if and only if $\fL \otimes \fO_{X_i}$ is ample on $X_i$, for each irreducible component $X_i$ of $X$.
  \item[(d)] Let $f : X \to Y$ be a finite surjective morphism, and let $\fL$ be an invertible sheaf on $Y$.
  Then $\fL$ is ample on $Y$ if and only if $f^* \fL$ is ample on $X$.
\end{itemize}

\begin{proof} $ $ \vspace{0pt}
\begin{itemize} [leftmargin=0cm]
  \item[(a)] If $\fF$ is any coherent sheaf on $Y$, then $f_*\fF$ is a coherent $\fO_X$-module (II, 5.5c)

  \item[(b)]

  \item[(c)]

  \item[(d)] By (b) and (c), we reduce to the case $X$ and $Y$ are irreducible, reduced, and proper schemes  over $A$.


\end{itemize} 
\end{proof}

\item[\textbf{10.}] Let $X$ be a projective scheme over a Noetherian ring $A$, and let $\fF^1 \to \fF^2 \to \cdots \to \fF^r$ be an exact sequence of coherent sheaves on $X$. Show that there is an integer $n_0$, such that for all $n \geq n_0$, the sequence of global sections
\[ \begin{tikzcd}
{\Gamma(X, \fF^1(n))} \arrow[r] & {\Gamma(X, \fF^2(n))} \arrow[r] & \cdots \arrow[r] & {\Gamma(X, \fF^r(n))}
\end{tikzcd} \]
is exact.

\begin{proof}
  Denote $f_i : \fF^i : \to \fF^{i + 1}$ for each $i = 1, \dots, n - 1$. 
  For every $i = 2, \dots, n-1$, we have the following three short exact sequence
  \[ \begin{tikzcd}
    0 \arrow[r] & \im{f_{i - 1}} \arrow[r]  & \fF^i \arrow[r]       & \im{f_{i}} \arrow[r]   & 0, \\
    0 \arrow[r] & \ker{f_{i - 1}} \arrow[r] & \fF^{i - 1} \arrow[r] & \im{f_{i-1}} \arrow[r] & 0, \\
    0 \arrow[r] & \im{f_i} \arrow[r]        & \fF^{i + 1} \arrow[r]       & \coker{f_i} \arrow[r]  & 0.
  \end{tikzcd} \]
  Every sheaf above is coherent (II, 5.7).
  By (5.2), there exists an integer $n_i$ such that for each $i > 0$ and each $n \geq n_i$, the cohomology vanishes for every sheaf, i.e., there is an exact sequence of global sections
  \[ \begin{tikzcd}
    0 \arrow[r] & \gsecx{\im{f_{i - 1}}(n)} \arrow[r]  & \gsecx{\fF^i(n)} \arrow[r]       & \gsecx{\im{f_{i}}(n)} \arrow[r]   & 0 \\
    0 \arrow[r] & \gsecx{\ker{f_{i - 1}}(n)} \arrow[r] & \gsecx{\fF^{i - 1}(n)} \arrow[r] & \gsecx{\im{f_{i-1}}(n)} \arrow[r] & 0 \\
    0 \arrow[r] & \gsecx{\im{f_i}(n)} \arrow[r]        & \gsecx{\fF^{i + 1}(n)} \arrow[r]       & \gsecx{\coker{f_i}(n)} \arrow[r]  & 0
  \end{tikzcd} \]
  which implies the sequence
  \[ \begin{tikzcd}
    \gsecx{\fF^{i - 1}(n)} \arrow[r] & \gsecx{\fF^{i}(n)} \arrow[r] & \gsecx{\fF^{i + 1}(n)}
  \end{tikzcd} \]
  is exact.
  Take $n_0 = \max{n_i}$ to obtain the desired integer.
\end{proof}

\end{enumerate}
\end{document}

% \item Recall from (II, Ex. 6.10) the definition of the Grothendieck group $K(X)$ of a Noetherian scheme $X$.
% \begin{itemize}
%   \item[(a)] Let $X$ be a projective scheme over a field $k$, and let $\fO_X(1)$ be a very ample invertible sheaf on $X$.
% Show that there is a (unique) additive homomorphism
%   \begin{equation*}
%     P : K(X) \to \Q[z]
%   \end{equation*}
%   such that for each coherent sheaf $\fF$ on $X$, $P(\gamma(\fF))$ is the Hilbert polynomial of $\fF$ (Ex. 5.2).
%   \item[(b)] Now let $X = \PP_k^r$. For each $i = 0, \dots, r$, let $L_i$ be a linear space of dimension $i$ in $X$. Then show that
%   \begin{itemize}
%     \item[(1)] $K(X)$ is the free Abelian group generated by $\{ \gamma(\fO_{K_i}) \mid i = 0, \dots, r \}$, and
%     \item[(2)] the map $P : K(X) \to \Q[z]$ is injective.
%   \end{itemize}
% \end{itemize}

% \item Let $k$ be a field, let $X = \PP_k^r$, and let $Y$ be a closed subscheme of dimension $q \geq 1$, which is a complete intersection (II, Ex. 8.4). Then:
% \begin{itemize}
%   \item[(a)] for all $n \in \Z$, the natural map
%   \begin{equation*}
%     H^0(X, \fO_X(n)) \to H^0(Y, \fO_Y(n))
%   \end{equation*}
%   is surjective. (This gives a generalization and another proof of (II, Ex. 8.4c), where we assumed $Y$ was normal.)
%   \item[(b)] $Y$ is connected;
%   \item[(c)] $H^i(Y, \fO_Y(n)) = 0$ for $0 < i < q$ and all $n \in \Z$;
%   \item[(d)] $p_a(Y) = \dim_k{H^q(Y, \fO_Y)}$.
% \end{itemize}

% \item Let $X$ (respectively, $Y$) be proper schemes over a Noetherian ring $A$. We denote by $\mathscr{L}$ an invertible sheaf.
% \begin{itemize}
%   \item[(a)] If $\mathscr{L}$ is ample on $X$, and $Y$ is any closed subscheme of $X$, then $i^* \mathscr{L}$ is ample on $Y$, where $i : Y \to X$ is the inclusion.
%   \item[(b)] $\mathscr{L}$ is ample on $X$ if and only if $\mathscr{L}_\text{red} = \mathscr{L} \otimes \fO_{X_\text{red}}$ is ample on $X$.
%   \item[(c)] Suppose $X$ is reduced. Then $\mathscr{L}$ is ample on $X$ if and only if $\mathscr{L} \otimes \fO_{X_i}$ is ample on $X_i$, for each irreducible component $X_i$ of $X$.
%   \item[(d)] Let $f : X \to Y$ be a finite surjective morphism, and let $\fL$ be an invertible sheaf on $Y$. Then $\fL$ is ample on $Y$ if and only if $f^* \mathscr{L}$ is ample on $X$.
% \end{itemize}

% \item \textit{A Nonprojective scheme.} We show the result of (Ex. 5.8) is false in dimension $2$. Let $k$ be an algebraically closed field of characteristic $0$, and let $X = \PP_k^2$. Let $\omega$ be the sheaf of differential $2$-forms (II, \S 8). Define an infinitesimal extension $X'$ of $X$ by $\omega$ by giving the element $\xi \in H^1(X, \omega \otimes \mathscr{T})$ defined as follows (Ex. 4.10). Let $x_0, x_1, x_2$ be the homogenous coordinates of $X$, let $U_0, U_2, U_2$ be the standard open covering, and let $\xi_{ij} = (x_j / x_i) d(x_i / x_j)$. This gives a Čech $1$-cocycle with values in $\Omega_X^1$, and since $\dim{X} = 2$, we have $\omega \otimes \mathscr{T} \cong \Omega^1$ (II, Ex. 5.16b). Now use the exact sequence
% \begin{equation*}
%   \cdots \to H^1(X, \omega) \to \pic{X'} \to \pic{X} \xrightarrow{\delta} H^2(X, \omega) \to \cdots
% \end{equation*}
% of (Ex. 4.6) and show $\delta$ is injective. We have $\omega \cong \fO_X(-3)$ by (II, 8.20.1), so $H^2(X, \omega) \cong k$. Since $\text{char}~k = 0$, you need only show that $\delta(\fO(1)) \neq 0$, which can be done by calculating in Čech cohomology. Since $H^1(X, \omega) = 0$, we see that $\pic{X'} = 0$. In particular, $X'$ has no ample invertible sheaves, so it is not projective.

% \textit{Note.} In fact, this result can be generalized to show that for any nonsingular projective surface $X$ over an algebraically closed field $k$ of characteristic $0$, there is an infinitesimal extension $X'$ of $X$ by $\omega$, such that $X'$ is not projective over $k$. Indeed, let $D$ be an ample divisor on $X$. Then $D$ determines an element $c_1(D) \in H^1(X, \Omega^1)$ which we use to define $X'$, as above. Then for any divisor $E$ on $X$ one can show that $\delta(\fL(E)) = (D.E)$, where $(D.E)$ is the intersection number (Chapter V), considered as an element of $k$. Hence, if $E$ is ample, $\delta(\mathscr{L}(E)) \neq 0$. Therefore, $X'$ has no ample divisors.

% On the other hand, over a field of characteristic $p > 0$, a proper scheme $X$ is projective if and only if $X_\text{red}$ is!
